\documentclass[11pt]{article}
\usepackage[margin=1in]{geometry}
\usepackage{amsmath,amssymb,amsthm}
\usepackage{booktabs}
\usepackage{array}
\usepackage[hidelinks]{hyperref}

\newtheorem{theorem}{Theorem}
\newtheorem{lemma}[theorem]{Lemma}
\newtheorem{corollary}[theorem]{Corollary}
\theoremstyle{definition}
\newtheorem{definition}[theorem]{Definition}
\theoremstyle{remark}
\newtheorem{remark}[theorem]{Remark}

\newcommand{\lean}[1]{\texttt{#1}}
\newcommand{\Qb}{Q}
\newcommand{\Eb}{E}

\title{Joint vertex covers and the envelope of Gowers's inverse induction}
\author{Formalization notes, Combinatorics/Ramsey (2026-10-11)}
\date{}

\begin{document}
\maketitle

\begin{abstract}
In Gowers's proof of Szemer\'edi's theorem, the structural step of Section~16
(the single-piece lift) must produce one multilinear map that agrees with the
\emph{remainder} $\varphi''$ on a dense piece of a box. The remainder is a
signed sum over the $2^k-1$ non-top vertices of the cube $\{0,1\}^k$, and each
vertex contributes a function that lives on a lower-dimensional face. If the
vertices are handled one after another, each at the density left by the
previous one, the degrees of the density losses multiply, and the resulting
envelope grows like $3^k$ in the doubly-logarithmic scale. That outgrows the
printed budgets $2^{d+10}$ (Theorem~18.1) and $2^{L+9}$ (Theorem~18.2) at
degree $10$ and length $14$. We show that the vertex \emph{covers} can be
composed at a precision fixed in advance, with a single pigeonhole at the end.
The losses then add, the envelope grows like $2^k$, and Theorems~18.1, 18.2
and Corollary~18.7 hold in every degree and length, at the printed constants
and with some room to spare. Every statement below is formalized in Lean
(module names are given in \S\ref{sec:lean}).
\end{abstract}

\section{Covers}

Throughout, $N$ is prime, a \emph{box} is a product of arithmetic progressions
in $\mathbb{Z}_N$ with a common difference, and $w(P)$ is the least side
length of a box $P$. A map is \emph{multilinear} in the sense of the paper.

\begin{definition}[Cover]
Let $\Gamma\subseteq\mathbb{Z}_N^d\times\mathbb{Z}_N$ be a relation and
$\Qb,\Eb\colon(0,1]\to\mathbb{R}$. We say $\Gamma$ is \emph{covered with controls
$(\Qb,\Eb)$} if for every $\theta\in(0,1]$ and every proper box $P$ there are
\begin{itemize}
\item a partition of $P$ into proper boxes $Q_1,\dots,Q_M$ with
  $w(Q_j)\ge w(P)^{\Eb(\theta)}$,
\item a set $H\subseteq P$ with $|H|\ge(1-\theta)|P|$,
\item for each $j$ at most $\Qb(\theta)$ multilinear maps
  $\mu_{j,1},\dots,\mu_{j,q}$,
\end{itemize}
such that every $(x,y)\in\Gamma$ with $x\in Q_j\cap H$ satisfies
$y=\mu_{j,i}(x)$ for some $i$.
\end{definition}

The precision $\theta$ is chosen by the user of the cover \emph{before} it is
applied, and the loss $\theta|P|$ is additive. A \emph{provider}, by contrast,
takes a set of density $t$ and returns a piece of density $c(t)$; composing
providers composes the functions $c$.

\section{Joint covers}

\begin{lemma}[Joint cover]\label{lem:joint}
Let $E$ be a finite index set, and for $e\in E$ let $f_e$ be a function on a
domain $D_e\subseteq\mathbb{Z}_N^d$ whose graph is covered with controls
$(\Qb_e,\Eb_e)$, with $\Qb_e,\Eb_e\ge0$. Fix precisions $\theta_e\in(0,1]$.
Then every proper box $P$ has a partition into proper boxes $R_1,\dots,R_M$, a
set $G\subseteq P$ and, for each $j$, at most $q\le\prod_e\Qb_e(\theta_e)$
tuples $(\nu_{j,i,e})_{e\in E}$ of multilinear maps, such that
\[
  |G|\ge\Bigl(1-\sum_{e}\theta_e\Bigr)|P|,\qquad
  w(R_j)\ge w(P)^{\prod_e\Eb_e(\theta_e)},
\]
and every $x\in R_j\cap G$ lying in all domains $D_e$ has an index $i$ with
$f_e(x)=\nu_{j,i,e}(x)$ for all $e$.
\end{lemma}

\begin{proof}
Induction on $E$. For $E=\emptyset$ take the trivial partition, $G=P$ and one
empty tuple. For $E\cup\{a\}$, apply the induction hypothesis to $P$, obtaining
cells $Q_j$, a good set $G$ and tuples $\nu_{j,i}$. Apply the cover of $f_a$ at
precision $\theta_a$ to every cell $Q_j$; this gives sub-cells $R_{j,l}$, good
sets whose union $G'$ satisfies $|G'|\ge(1-\theta_a)|P|$ (the cells partition
$P$), and maps $\mu_{j,l,i'}$. Take the sub-cells $R_{j,l}$ as the new
partition, the good set $G\cap G'$, of size at least
$|G|+|G'|-|P|\ge(1-\sum_E\theta_e-\theta_a)|P|$, and the product tuples
$(\nu_{j,i},\mu_{j,l,i'})$. The width bound follows from
$w(R_{j,l})\ge w(Q_j)^{\Eb_a}\ge w(P)^{\Eb_a\prod_E\Eb_e}$.
\end{proof}

\begin{corollary}[One piece]\label{cor:piece}
In the setting of Lemma~\ref{lem:joint}, put $s=\sum_e\theta_e$ and
$Q=\prod_e\Qb_e(\theta_e)$. If $H\subseteq P$ lies in every domain $D_e$ and
$|H|\ge\theta|P|$ with $\theta>s$, then some cell $R$ and some maps $\nu_e$
satisfy
\[
  \bigl|\{x\in H\cap R: f_e(x)=\nu_e(x)\text{ for all }e\}\bigr|
  \ge\frac{\theta-s}{Q}\,|R| .
\]
\end{corollary}

\begin{proof}
$|H\cap G|\ge(\theta-s)|P|$, so some cell $R$ has
$|H\cap G\cap R|\ge(\theta-s)|R|$. Each of its points is covered by one of at
most $Q$ tuples; the most popular tuple covers a $1/Q$ fraction.
\end{proof}

The only density loss is the single factor $1/Q$, and $Q$ is a product of
counts evaluated at \emph{fixed} precisions. If $\Qb_e(\theta)\le(C/\theta)^{D_e}$,
the density has degree $1+\sum_e D_e$ in $1/\theta$: the degrees add.

\section{The remainder of the single-piece step}

On the good domain of the step, the remainder is
$\varphi''(x)=-\sum_{e\ne\mathbf 1}(-1)^{k+|e|}\varphi_e(x)$, where $\varphi_e$ is
the translated cube vertex $e$, a pullback of $\varphi$ to the face spanned by
the directions $S_e$, of dimension $|e|+1$. By induction on the dimension, the
face carries a cover of degree $D_{|S_e|}$. It lifts to the ambient box without
any loss in the precision: the count becomes $\max(\Qb,3^{k+1})$ and the
exponent $\Eb/16^{k+1-|S_e|}$. Applying Corollary~\ref{cor:piece} with
$\theta_e=\theta_2/2^{k+1}$ to the dense good set (density $\theta_2$) gives one
map $M''=-\sum_e(-1)^{k+|e|}\nu_e$ on a piece of density
\[
  \rho_1=\frac{\theta_2/2}{\prod_{e\ne\mathbf 1}\max\bigl(\Qb_{|S_e|}(\theta_2/2^{k+1}),3^{k+1}\bigr)}
\]
and width at least $N^{\prod_e\Eb_e/16^{k+1-|S_e|}}$.

\begin{remark}[Nesting versus joint covers]
With the class degrees $D_l=2^{68\cdot2^l}$, nesting one provider per vertex
gives $\log_2\deg\rho_1^{-1}\approx\sum_{e}68\cdot2^{|e|+1}=136\cdot3^k$. The
joint cover gives $\log_2\bigl(1+\sum_eD_{|S_e|}\bigr)\le68\cdot2^k+k$.
\end{remark}

\section{The envelope}

Write $X=2/\alpha$. The inverse induction carries bounds
$\beta^{-1},\sigma^{-1}\le X^{2^{\mu}}$ from degree $d-1$ to degree $d$. With the
joint-cover remainder, one step in dimension $k$ costs
\[
  \mu_{\mathrm{new}}=\mu+\underbrace{(68\cdot2^k+\mu_0+k+3)}_{\rho_1}
  +2\cdot68\cdot2^k+O(k^2),
\]
where $\mu_0$ bounds the base budget of the cover controls. With the standard
$\mu_0\approx64\cdot2^k$ one step costs $268\cdot2^k+O(k^2)$. Hence

\begin{theorem}
For every $k\ge1$ the Theorem~18.1 envelope satisfies $\mu_k\le1000\cdot2^k$,
and the Theorem~18.2 envelope (which also carries thresholds) satisfies
$\mu_k\le1000\cdot2^k$ as well. Both are well inside the printed budgets
$2^{d+10}=4096\cdot2^{d-2}$ and $2^{L+9}=8192\cdot2^{L-4}$. Consequently
Theorems~18.1 and 18.2 and Corollary~18.7 hold in every degree and length.
\end{theorem}

\section{Sharper constants}

The same proofs give smaller tower constants. Moreover, the class base
$V=1+\sum_{j<k}(K_j+p_j)$ satisfies $V\le2^{2^{2^{k+4}+1}}$, so the base budget
$\mu_0=2^{k+5}+9$ suffices; one step then costs $236\cdot2^k+O(k^2)$ and the
envelopes grow like $472\cdot2^k<512\cdot2^k$.

\begin{center}\small
\begin{tabular}{lccc}
\toprule
& printed & standard base & tight base\\
\midrule
Theorem 18.1, exponent $\alpha^{2^{2^{d+c}}}$ & $c=10$ & $c=8$ ($d\ge4$) & $c=7$ (every $d$)\\
Theorem 18.2, threshold with $L+c$ & $c=9$ & $c=6$ (every $L$) & $c=5$ (every $L$)\\
Corollary 18.7, threshold with $L+c$ & $c=9$ & $c=6$ (every $L$) & $c=5$ (every $L$)\\
$W(k,r)\le2^{2^{r^{2^{2^{k+c}}}}}$, $r\ge2$ & $c=9$ & --- & $c=5$\\
Theorem 1.3, density $(\log\log N)^{-2^{-2^{k+c}}}$ & $c=9$ & --- & $c=5$\\
\bottomrule
\end{tabular}
\end{center}

The $W(k,r)$ row is Theorem~18.2 at density $\delta=1/r$: in an $r$-colouring of
$[1,N]$ one colour class has at least $N/r$ elements. This gives Gowers's
van der Waerden bound, sharpened from $k+9$ to $k+5$. The headline
Theorem~1.3 follows from Theorem~18.2 by a calculation in which the constant
enters only through $2^{k+9}$; repeated with $k+5$, it gives the density
exponent $2^{-2^{k+5}}$.

The method stops at these constants. The envelope recursion grows like
$472\cdot2^k$, and $c=6$ for Theorem~18.1 (or $c=4$ for Theorem~18.2) would
need a step cost below $128\cdot2^k$, while the vertex covers and the two slice
controls alone cost $204\cdot2^k$.

For Theorem~18.2 at $c=5$ one more point matters. The modulus threshold of
the inverse step contains two pure constants, the Section~5 refinement
constant $\le2^{2^{40d^3+2}}$ and a boundary constant $\le2^{2^{1083}}$. A
cruder bound charges them at the level of powers of $X$, which puts a cubic
term $40(k+2)^3$ into the envelope and blocks $L+5$ for $L\le10$. The
threshold is only needed below $\exp(\exp(X^{2^\mu}))$, and
$2^{2^m}\le\exp(X^{2^j})$ as soon as $m\le2^j$. At that level the
constants cost only $k+16$, and $L+5$ holds in every length.

\section{Theorem 16.2 and Corollary 16.11 with polynomial controls}

The printed Theorem~16.2 restricts a relation with the product property to a
set of density $1-\theta$ on which it is multiply linear, with graph count
$q(r^{-1}\rho,\gamma,k)^r$ at precision $\rho$, where
$r=\gamma^{-2}s(\theta,\gamma,k)$ and $s=(2/\theta\gamma)^{2^{2^{k+6}}}$. This
count is exponential in a polynomial of $1/\theta$. The covers built for the
dimension induction give a count that is \emph{polynomial} in $1/(\gamma\theta\rho)$:
\[
  Q(\gamma,\theta,\rho)=\max\Bigl(1,\Bigl(\frac{4V}{\gamma\theta\rho}\Bigr)^{2^{68\cdot2^k}}\Bigr),
  \qquad V=2^{2^{2^{k+4}+1}},
\]
with width exponent $Q^{-1}$, in every dimension $k\ge1$
(\lean{theorem\_16\_2\_polynomial}). This form implies the printed one
(\lean{theorem\_16\_2\_polynomial\_implies\_printed}).

The same holds for Corollary~16.11. For a function that is not uniform of
degree $k+2$, the single-piece step with the joint-cover remainder gives a
box of dimension $k+1$, of width $N^e$, on which a multilinear frequency is
large on an $e$ fraction, with
\[
  e^{-1}\le(2/\alpha)^{2^{236\cdot2^k+k+24}}
\]
for every $k\ge2$ (\lean{corollary\_16\_11\_polynomial}). The printed exponent
is $(\alpha/8)\,c(r^{-1}\alpha/8,\alpha/2)^{r}$ with $r$ polynomial in
$1/\alpha$, so its reciprocal is exponential in a polynomial of $1/\alpha$.

\section{Further questions}

\begin{enumerate}
\item \emph{Joint slices.} The single-piece step applies its slice control
  twice in succession, $\theta\,c(c(\theta))$, and so pays $2\cdot68\cdot2^k$
  in the step cost. Composing the two slice covers jointly, as in
  Lemma~\ref{lem:joint}, should cost about $68\cdot2^k$. The envelope would
  drop from about $472\cdot2^k$ to about $336\cdot2^k$. This changes none of
  the constants in the table.
\item \emph{The class degree.} The face covers have degree
  $D_l=2^{68\cdot2^l}$, the fixed point of a class recursion that pays the
  Section~16 exponent $\theta_1=(\theta\gamma/2)^{2^{2^{k+5}}}$ (about
  $32\cdot2^k$ in the log-degree) once per dimension. Is there a class
  construction with log-degree near one $\theta_1$ per dimension?
\item \emph{The printed $\theta_1$.} Even with joint slices and such a class
  degree, a step would cost about $134\cdot2^k$, above the $128\cdot2^k$ that
  $c=6$ for Theorem~18.1 or $c=4$ for Theorem~18.2 would need. Can the
  double-exponential exponent $2^{2^{k+5}}$ of $\theta_1$ in Lemmas 16.5 and
  16.7 be lowered? Only that would move these constants further. The
  exponent is created by the random selection of Lemma~15.5, whose Riesz
  kernel $(1+\cos)/2$ gives a respected arrangement of $M=2^{k+4}$ vertices
  an advantage of only $1+2^{1-M}$ per stage. A single Fej\'er stage of
  order $m\approx M/(\alpha\eta)$ gives an advantage of about
  $1/(\alpha\eta)$ at density cost $(\alpha\eta/2M)^{M}$, provided Lemma~15.4
  is extended from ternary to $[-m,m]$ coefficients. Then $\theta_1$ has
  exponent about $(k+7)\,2^{4k+10}$. An exact model of the class recursion
  gives a class degree $2^{14\cdot2^l}$ and step cost below $128\cdot2^k$:
  Theorem~18.1 at $c=6$ and Theorem~18.2 at $c=4$ (notes, L.16; not
  formalized).
\end{enumerate}

\section{Formalization}\label{sec:lean}

All statements are in the directory
\lean{Combinatorics/Ramsey/Lean/GowersSzemeredi}.

\begin{center}\small
\begin{tabular}{>{\raggedright\arraybackslash}p{0.28\textwidth}>{\raggedright\arraybackslash}p{0.64\textwidth}}
\toprule
Statement & Lean\\
\midrule
Lemma~\ref{lem:joint}, Corollary~\ref{cor:piece} & \lean{joint\_cover\_on}, \lean{JointCoverOn.piece} (\lean{Proofs16JointCovers})\\
the remainder piece & \lean{remainder\_piece\_of\_vertex\_covers}\\
the single-piece step & \lean{single\_piece\_lift\_global\_cover} (\lean{Proofs16SinglePieceCoverGlobal})\\
$\rho_1$ bound & \lean{rho1\_inv\_leC} (\lean{Proofs18CoverDensity})\\
envelopes & \lean{polyGenMuC\_le}, \lean{genMuC\_le}, \lean{genMuC\_le\_sharp}\\
Theorems 18.1, 18.2, Cor.~18.7 & \lean{theorem\_18\_1\_explicit}, \lean{theorem\_18\_2\_explicit},\\
& \lean{corollary\_18\_7\_explicit} (\lean{Proofs18CoverExplicit})\\
sharper constants & \lean{Proofs18CoverSharper}, \lean{Proofs18CoverTightBase},\\
& \lean{Proofs18CoverColorings}, \lean{Proofs18CoverSharperExplicit},\\
& \lean{Proofs18CoverSharperLowDegrees}, \lean{Proofs18CoverSharperShortLengths},\\
& \lean{Proofs18CoverThresholdsSharp}, \lean{Proofs18CoverVanDerWaerden},\\
& \lean{Proofs18CoverHeadline},\\
& \lean{Proofs18CoverSharperAllDegrees}\\
polynomial Theorem 16.2 & \lean{Proofs16PolynomialStructure}, \lean{theorem\_16\_2\_polynomial}\\
polynomial Cor.~16.11 & \lean{Proofs16PolynomialCorollary}, \lean{corollary\_16\_11\_polynomial}\\
\bottomrule
\end{tabular}
\end{center}

Every result above is kernel-checked. The production audit of 2026-10-11,
\begin{center}\small
\lean{check\_gowers.py GowersSzemeredi.CompletionAudit},
\end{center}
built 1,440 modules from source with Lean~4.32.0, checked the exact types of
all 120 numbered catalogue statements and of the improved-bound theorems,
and found only the axioms \lean{propext}, \lean{Classical.choice} and
\lean{Quot.sound}. The extensions to every degree ($c=7$) and to every
length ($c=6$ and $c=5$), the van der Waerden and Theorem~1.3 rows, and the
polynomial Theorem~16.2 and Corollary~16.11 postdate that audit. They are so far checked against
faithful stubs, and \lean{CompletionAudit.lean} now includes them. The
audit's record is
\begin{center}\small
\lean{Research/GowersSzemeredi/verification/}\\
\lean{2026-10-11-numbered-and-bounds.json}.
\end{center}
The research notes
(\lean{SECTION16\_RESEARCH\_NOTES.md}, \S\S L.11--L.16) record the development.

\end{document}
